\section{A conditional application to a heteroclinic parameter}
\label{sec:heteroclinic-connection}

The local matching function in \cref{thm:rw-local-fold-matching} compares
two finite-interval traces.  A heteroclinic parameter requires, in addition,
a defining function for global invariant-manifold membership.  We give
sufficient hypotheses for transferring the local sensitivity to such a
function.  The global invariant-manifold construction and the comparison
estimates stated below are not established in this paper.

\subsection{A fixed recovery law with two outer equilibria}
\label{subsec:global-modification}

At an equilibrium, the original recovery equation
\eqref{eq:rw-original-recovery} forces
\(\pi_N^T(v-\mathbf1)=\delta^2\nu\).  To prescribe two equilibria outside
the fold neighborhood, we change the recovery law while keeping it fixed
near the fold.  Write
\[
 \mathcal N_N(x)=c_N\circ x^{\circ2}+\frac\beta3x^{\circ3}.
\]
The uniform inverse bound for \(A_N\) and the product estimate
\eqref{eq:fold-product-bound} give a radius \(r_0>0\), independent of
\(N\), and a local critical curve determined by
\begin{equation}
 A_Nz_{0,N}(r)=P_{\perp,N}\mathcal N_N
       (r\mathbf1+z_{0,N}(r)),
 \qquad z_{0,N}(r)\in E_N,
 \qquad \norm{z_{0,N}(r)}_N\le Cr^2.
 \label{eq:rw-critical-curve-equation}
\end{equation}
Indeed, applying \(A_N^{-1}\) to the right-hand side is a contraction on a
ball of radius \(Cr^2\) for \(\abs r\le r_0\), with constants independent
of \(N\).  Set
\begin{equation}
 \begin{aligned}
 V_{0,N}(r)&=\mathbf1+r\mathbf1+z_{0,N}(r),\\
 w_{0,N}(r)&=\frac23-\pi_N^T\mathcal N_N
                       (r\mathbf1+z_{0,N}(r)).
 \end{aligned}
 \label{eq:rw-critical-curve}
\end{equation}
Every delayed difference vanishes on a constant history, so this curve
consists of equilibria of the frozen voltage equation for every \(\eta\).

Fix \(0<r_{\rm loc}<r_A<r_0\) and one function
\(g\in C^{12}(\mathbb R)\) satisfying
\begin{equation}
 \begin{aligned}
 g(r)&=1 \quad(\abs r\le r_{\rm loc}),
 &g(r)&>0 \quad(\abs r<r_A),\\
 g(\pm r_A)&=0,
 &g'(-r_A)&>0>g'(r_A),\\
 \norm{g}_{C^{12}([-r_0,r_0])}&\le G_*.
 \end{aligned}
 \label{eq:rw-global-recovery-multiplier}
\end{equation}
Such a function can be obtained by a smooth interpolation between the
constant function near zero and \(1-r^2/r_A^2\) near \(\pm r_A\).
It is chosen independently of \(N,\delta,\nu,\eta\).
For \(r(v_t)=\pi_N^T\{v(t)-\mathbf1\}\), replace the recovery equation by
\begin{equation}
 \dot w(t)=\delta^2\bigl(r(v_t)-\delta^2\nu\bigr)g(r(v_t)).
 \label{eq:rw-modified-recovery}
\end{equation}
This equation and \eqref{eq:rw-voltage-rfde} agree with the original RFDE
on \(\abs{r(v_t)}\le r_{\rm loc}\) and retain the projection identity
\eqref{eq:rw-projection-identity}.  They have the exact equilibria
\begin{equation}
 Z_N^\pm=
 \left(\theta\longmapsto V_{0,N}(\pm r_A),
       w_{0,N}(\pm r_A)\right)
 \in\mathcal H_{N,\delta}.
 \label{eq:rw-outer-equilibria}
\end{equation}

\subsection{Comparison with a global defining function}

When \(Z_N^+\) is hyperbolic with a one-dimensional unstable manifold, let
\(W_{\rm in}^u(Z_N^+)\) denote the branch on which \(r\) initially
decreases from \(r_A\).  We seek a connection from this branch to
\(Z_N^-\), with convergence in \(\mathcal H_{N,\delta}\) at both ends.
A scalar defining function may be obtained on a common transverse section
as follows: if the incoming branch meets the section at
\((\psi_A,u_A)\) and the stable-manifold section is the graph
\(u=F^-(\psi)\) on a domain containing \(\psi_A\), then
\(G=u_A-F^-(\psi_A)\).  This construction requires an actual
stable-manifold section and its common domain; a finite terminal condition
alone does not provide it.

\begin{corollary}[Conditional heteroclinic connection]
\label[corollary]{cor:rw-conditional-connection}
Assume the hypotheses of \cref{thm:rw-local-fold-matching} and fix the
recovery law \eqref{eq:rw-modified-recovery}.  Choose \(c_0>0\) and
\(0<\eta_0<\eta_*\), independently of \(N\), such that
\([\nu_{0,N}-c_0,\nu_{0,N}+c_0]\subset I_\nu\) for every \(N\).
For every sufficiently small \(\delta>0\), set
\[
 \mathcal P_{N,\delta}
 =\{(\nu,\eta):\abs{\nu-\nu_{0,N}}\le c_0,
                    \ \norm{\eta}_{\mathcal T_N}\le\eta_0\}.
\]
Suppose the following hold for every \(N\) on this parameter set.
The equilibria \(Z_N^\pm\) are hyperbolic, with
\[
 \dim W^u(Z_N^+)=1,
 \qquad \operatorname{codim}W^s(Z_N^-)=1,
\]
and the specified incoming branch and stable-manifold section form \(C^1\)
parameter families in a common transverse section of compatible histories.
There is a scalar function \(G_{N,\delta}^g\), of class \(C^1\) on a
neighborhood of \(\mathcal P_{N,\delta}\), such that
\begin{equation}
 G_{N,\delta}^g(\nu,\eta)=0
 \quad\Longleftrightarrow\quad
 W_{\rm in}^u(Z_N^+)\subset W^s(Z_N^-).
 \label{eq:connection-zero-set}
\end{equation}
With \(E_{N,\delta}=G_{N,\delta}^g-D_N^{\rm fin}\), assume the comparison
bounds, with \(C\) independent of \(N\) and \(\delta\),
\begin{equation}
 \abs{E_{N,\delta}}+\abs{\partial_\nu E_{N,\delta}}
       \le C\delta^2,
 \qquad
 \norm{D_\eta E_{N,\delta}}_{\mathcal T_N^*}
       \le C\delta^3.
 \label{eq:rw-conditional-connection-hypothesis}
\end{equation}
Then, after reducing the upper bound on \(\delta\) independently of
\(N\), for every \(\norm{\eta}_{\mathcal T_N}<\eta_0\) there is a unique
parameter
\begin{equation}
 \nu_{c,N}^g(\delta,\eta)
       \in(\nu_{0,N}-c_0,\nu_{0,N}+c_0),
 \qquad
 \mu_{c,N}^g(\delta,\eta)=\delta^2\nu_{c,N}^g(\delta,\eta),
 \label{eq:rw-exact-root}
\end{equation}
at which the specified branch is a heteroclinic orbit from \(Z_N^+\) to
\(Z_N^-\).  This orbit is unique modulo time translation on that branch,
and \(\abs{\nu_{c,N}^g-\nu_{0,N}}\le C\delta\).
\end{corollary}

\begin{proof}
Put \(a_0=\sqrt{2\pi}\).  The level and slope estimates of
\cref{thm:rw-local-fold-matching} and
\eqref{eq:rw-conditional-connection-hypothesis} give
\begin{equation}
 \delta^{-1}G_{N,\delta}^g
       =a_0(\nu-\nu_{0,N})+O(\delta),
 \qquad
 \partial_\nu G_{N,\delta}^g=\delta a_0+O(\delta^2).
 \label{eq:connection-function-estimates}
\end{equation}
For sufficiently small \(\delta\), uniformly in \(N\) and \(\eta\),
the first expression has opposite signs at
\(\nu=\nu_{0,N}\pm c_0\), while
\(\partial_\nu G_{N,\delta}^g\ge\delta a_0/2>0\).
The intermediate value theorem and strict monotonicity give exactly one
root.  Evaluating the first estimate at this root gives its distance
\(O(\delta)\) from \(\nu_{0,N}\).
By \eqref{eq:connection-zero-set}, the incoming history then belongs to
\(W^s(Z_N^-)\); it already lies on the specified unstable branch of
\(Z_N^+\).  Its complete orbit has the required limits.  A branch of a
one-dimensional unstable manifold consists of one orbit modulo time
translation, giving the stated branchwise uniqueness.
\end{proof}

\begin{corollary}[Sensitivity of the conditional connection parameter]
\label[corollary]{cor:rw-root-sensitivity}
Under the hypotheses of \cref{cor:rw-conditional-connection}, the function
\(\eta\mapsto\mu_{c,N}^g(\delta,\eta)\) is \(C^1\) for fixed
\(\delta>0\), and
\begin{align}
 \norm{D_\eta\mu_{c,N}^g(\delta,\eta)-\delta^3\Lambda_N}
 &\le C\{\delta^4+\delta^3\norm{\eta}_{\mathcal T_N}\},
 \label{eq:rw-root-derivative}\\
 \abs{\mu_{c,N}^g(\delta,\eta)-\mu_{c,N}^g(\delta,0)
                       -\delta^3\Lambda_N(\eta)}
 &\le C\{\delta^4\norm{\eta}_{\mathcal T_N}
                +\delta^3\norm{\eta}_{\mathcal T_N}^2\}.
 \label{eq:rw-root-increment}
\end{align}
The derivative norm is that of \(\mathcal T_N^*\), and the constants are
independent of \(N\).
\end{corollary}

\begin{proof}
The positive \(\nu\)-derivative in
\eqref{eq:connection-function-estimates} permits implicit differentiation;
uniqueness patches the resulting local \(C^1\) root functions over the
\(\eta\)-ball.  The structural derivative estimate in
\cref{thm:rw-local-fold-matching} and the second comparison bound give,
on the root graph,
\[
 D_\eta G_{N,\delta}^g
 =-\delta^2a_0\Lambda_N+R_\eta,
 \qquad
 \norm{R_\eta}\le C\{\delta^3+\delta^2\norm{\eta}_{\mathcal T_N}\}.
\]
Writing \(\partial_\nu G_{N,\delta}^g=\delta a_0+R_\nu\), where
\(\abs{R_\nu}\le C\delta^2\), yields
\[
 D_\eta\nu_{c,N}^g-\delta\Lambda_N
 =\frac{-R_\eta-\delta\Lambda_NR_\nu}
             {\delta a_0+R_\nu}.
\]
The uniform bound for \(\Lambda_N\) in
\eqref{eq:rw-functional-uniform-bound} therefore gives
\(\norm{D_\eta\nu_{c,N}^g-\delta\Lambda_N}
 \le C(\delta^2+\delta\norm{\eta}_{\mathcal T_N})\).
Multiplication by \(\delta^2\) proves
\eqref{eq:rw-root-derivative}.  Integrating this estimate along the segment
\(t\eta\), \(0\le t\le1\), proves \eqref{eq:rw-root-increment}.
\end{proof}

\begin{remark}[Centered conormal]
For the centered and rescaled connection locus, put
\begin{equation}
 \Xi_{N,\delta}^g(\eta)
 =\delta^{-3}\{\mu_{c,N}^g(\delta,\eta)
                       -\mu_{c,N}^g(\delta,0)\}.
 \label{eq:rw-centered-root-graph}
\end{equation}
At \(\eta=0\), its conormal normalized to have first component one obeys
\begin{equation}
 \norm{(1,-D_\eta\Xi_{N,\delta}^g(0))-(1,-\Lambda_N)}
       \le C\delta.
 \label{eq:rw-conormal-limit}
\end{equation}
This follows directly from \eqref{eq:rw-root-derivative}.  Multiplying a
defining function by a nowhere-zero \(C^1\) factor does not change its
conormal line at a zero.  Two fixed recovery laws satisfying the hypotheses
may have different values of \(\mu_{c,N}^g(\delta,0)\); each separately
centered locus has the same leading conormal in these parameter coordinates.
\end{remark}
